\section{Evaluation}
\label{sec:evaluation}

All results are from live PLC runs. ASR values are lower bounds under the safety shield: an unconstrained adversary could achieve higher ASR. This constraint applies uniformly and does not affect within-table comparisons.

\subsection{RQ1: Context Ablation}
\label{sec:eval:rq1}

Table~\ref{tab:context_ablation} shows how context depth affects attack capability (Qwen2.5-3B, no defense).

% Table 3 --- RQ1 Context Ablation
% Generated by: python -m cpsforge.analysis.generate_tables
% RQ-specific filtering: limited to 3 runs per (scene, context) cell
\begin{table}[t]
\centering
\caption{RQ1: Context ablation for online MITM attacker (base Qwen2.5-3B; no RQ4 defense active; safety shield always active).}
\label{tab:context_ablation}
\small
\setlength{\tabcolsep}{3pt}
\begin{tabular}{llrrrrrr}
\toprule
\textbf{Scene} & \textbf{Context} & \textbf{Runs} & \textbf{Prop.} & \textbf{Writes} & \textbf{Succ.} & \textbf{ASR\%} & \textbf{Lat.(ms)} \\
\midrule
Level Control & Minimal & 3 & 23 & 2 & 0 & 0.0 & 13086 \\
 & Partial & 3 & 20 & 13 & 0 & 0.0 & 18429 \\
 & Full & 3 & 5 & 0 & 0 & 0.0 & 7863 \\
\midrule
Sort.~Weight & Minimal & 3 & 13 & 13 & 7 & 53.8 & 17929 \\
 & Partial & 3 & 12 & 5 & 2 & 40.0 & 19789 \\
 & Full & 3 & 0 & 0 & 0 & 0.0 & 11180 \\
\midrule
Sort.~Height & Minimal & 3 & 30 & 30 & 1 & 3.3 & 16906 \\
 & Partial & 3 & 30 & 25 & 1 & 4.0 & 22944 \\
 & Full & 3 & 30 & 30 & 4 & 13.3 & 22263 \\
\bottomrule
\end{tabular}
\vspace{2pt}
\noindent\footnotesize{3 repeats per configuration; ASR computed post-hoc.}
\end{table}

\textbf{Finding 1: Self-deterrence under rich context.} Full context produces \emph{zero} attack proposals on Sorting by Weight---the model decides ``wait'' on every call across all repeats. On Level Control, full context reduces proposals from 23 (minimal) to 5, all shield-blocked. Rich operational context (phase inference, safety rules) makes the 3B model overly cautious.

\textbf{Finding 2: Self-deterrence is scene-dependent.} On Sorting by Height, full context is the \emph{most} effective tier (13.3\% ASR vs.\ 3.3\% at minimal)---the discrete state-machine logic benefits from richer phase context. On Sorting by Weight, partial context achieves 40\% ASR while full drops to zero proposals.

\textbf{Finding 3: Context changes strategy.} Minimal context produces brute-force overrides. Partial shifts toward setpoint manipulation. Full, when it does attack (Sort.~Height), targets a narrow set of tags with deliberate overrides. More context makes the model more selective but also more conservative.

\textbf{Implication:} Exposing operational context may itself serve as a defense---deterrence through transparency.

\subsection{RQ2: Fine-Tuning Impact}
\label{sec:eval:rq2}

Table~\ref{tab:finetune_comparison} compares base vs.\ QLoRA-finetuned Qwen2.5-3B.

% Table RQ2 --- Fine-Tuning Comparison
% Generated by: python -m cpsforge.analysis.generate_tables
% RQ-specific filtering: limited to 3 runs per (scene, context, model) cell
\begin{table}[t]
\centering
\caption{RQ2: Fine-tuning effect on attack capability (online MITM, no defense). Base = Qwen2.5-3B-Instruct; FT = same model with QLoRA v2 adapter. VAR = schema-valid outputs / LLM calls.}
\label{tab:finetune_comparison}
\small
\setlength{\tabcolsep}{3pt}
\resizebox{\columnwidth}{!}{%
\begin{tabular}{lllrrrrrrr}
\toprule
\textbf{Scene} & \textbf{Ctx} & \textbf{Model} & \textbf{Runs} & \textbf{LLM calls} & \textbf{VAR\%} & \textbf{Prop.} & \textbf{Writes} & \textbf{Succ.} & \textbf{ASR\%} \\
\midrule
Level Control & Minimal & Base & 3 & 27 & 100.0 & 23 & 2 & 0 & 0.0 \\
 & Minimal & FT & 3 & 30 & 100.0 & 30 & 30 & 11 & 36.7 \\
 & Full & Base & 3 & 15 & 100.0 & 5 & 0 & 0 & 0.0 \\
 & Full & FT & 3 & 34 & 0.0 & 0 & 0 & 0 & 0.0 \\
\midrule
Sort.~Weight & Minimal & Base & 3 & 14 & 92.9 & 13 & 13 & 7 & 53.8 \\
 & Minimal & FT & 3 & 30 & 100.0 & 30 & 16 & 0 & 0.0 \\
 & Full & Base & 3 & 5 & 40.0 & 0 & 0 & 0 & 0.0 \\
 & Full & FT & 3 & 60 & 0.0 & 0 & 0 & 0 & 0.0 \\
\midrule
Sort.~Height & Minimal & Base & 3 & 30 & 100.0 & 30 & 30 & 1 & 3.3 \\
 & Minimal & FT & 3 & 30 & 100.0 & 30 & 30 & 3 & 10.0 \\
 & Full & Base & 3 & 30 & 100.0 & 30 & 30 & 4 & 13.3 \\
 & Full & FT & 3 & 45 & 0.0 & 0 & 0 & 0 & 0.0 \\
\bottomrule
\end{tabular}%
}
\vspace{2pt}
\noindent\footnotesize{3 repeats per configuration; ASR computed post-hoc.}
\end{table}

\textbf{Finding 4: Fine-tuning enables attacks on Level Control.} The finetuned model achieves 36.7\% ASR (11/30 writes) on Level Control minimal, vs.\ 0\% for the base model. Training teaches the model to propose shield-legal values.

\textbf{Finding 5: Schema regression under full context.} The finetuned model achieves VAR=0\% under full context on all scenes---every output fails schema validation. This is a prompt-length distributional mismatch: training examples come from shorter prompts; full-context prompts ($\sim$1.5--2K tokens) fall outside the training distribution. Addressable by augmenting training data with full-context examples.

\textbf{Finding 6: Defense capability preserved.} The finetuned LLM defender retains 100\% blocking on Sort.~Height (Table~\ref{tab:defense_finetuned}), confirming dual-role training preserves defense while modifying attack behavior.

\subsection{RQ3: Online vs.\ Static Attack Mode}
\label{sec:eval:rq3}

Table~\ref{tab:attacker_comparison} compares three attacker modes (no defense).

% Table 4 --- RQ3 Attacker Comparison
% Generated by: python -m cpsforge.analysis.generate_tables
% RQ-specific filtering: limited to 3 runs per (scene, attacker) cell
\begin{table}[t]
\centering
\caption{RQ3: Attacker comparison (base Qwen2.5-3B, no defense). Random = bounded random perturbations; Static LLM = one-shot batch; Online MITM = core contribution.}
\label{tab:attacker_comparison}
\small
\setlength{\tabcolsep}{3pt}
\begin{tabular}{llrrrrr}
\toprule
\textbf{Scene} & \textbf{Attacker} & \textbf{Runs} & \textbf{Prop.} & \textbf{Writes} & \textbf{Succ.} & \textbf{ASR\%} \\
\midrule
Level Control & Random & 3 & 10 & 10 & 1 & 10.0 \\
 & Static LLM & 3 & 0 & 0 & 0 & 0.0 \\
 & Online MITM & 3 & 14 & 0 & 0 & 0.0 \\
\midrule
Sort.~Weight & Random & 3 & 0 & 0 & 0 & 0.0 \\
 & Static LLM & 3 & 0 & 0 & 0 & 0.0 \\
 & Online MITM & 3 & 13 & 13 & 7 & 53.8 \\
\midrule
Sort.~Height & Random & 3 & 30 & 30 & 4 & 13.3 \\
 & Static LLM & 3 & 0 & 0 & 0 & 0.0 \\
 & Online MITM & 3 & 30 & 30 & 1 & 3.3 \\
\bottomrule
\end{tabular}
\vspace{2pt}
\noindent\footnotesize{3 repeats per configuration; ASR computed post-hoc.}
\end{table}

\textbf{Finding 7: Static LLM generates zero valid proposals.} The one-shot batch LLM produces no schema-valid outputs---the 350-token budget is insufficient for a multi-object JSON array. This reflects our experimental design constraint, not a fundamental limitation of one-shot attacks.

\textbf{Finding 8: Online MITM is the only effective attacker on complex scenes.} On Sort.~Weight (35 tags), the online MITM achieves 53.8\% ASR while random and static generate zero proposals. The LLM's ability to reason about control logic is essential on complex scenes.

\textbf{Finding 9: Random baseline is competitive on discrete scenes.} On Sort.~Height, random achieves 13.3\% ASR vs.\ 3.3\% for the online MITM---binary actuators allow random overrides to accidentally align with vulnerable states.

\subsection{RQ4: Defense Effectiveness}
\label{sec:eval:rq4}

Table~\ref{tab:defense_comparison} (base model) and Table~\ref{tab:defense_finetuned} (finetuned) present defense results.

% Table 5 --- RQ4 Defense Comparison
% Generated by: python -m cpsforge.analysis.generate_tables
% RQ-specific filtering: limited to 3 runs per (scene, defense) cell
\begin{table}[t]
\centering
\caption{RQ4: Defense comparison for online MITM attacker (base Qwen2.5-3B). Prevention rate = fraction of attack proposals blocked before PLC write. ``None'' = no RQ4 defense; the mandatory safety shield (always active) accounts for any nonzero Blk.\ in ``None'' rows.}
\label{tab:defense_comparison}
\small
\setlength{\tabcolsep}{2pt}
\resizebox{\columnwidth}{!}{%
\begin{tabular}{llrrrrrrrrrr}
\toprule
\textbf{Scene} & \textbf{Defense} & \textbf{Runs} & \textbf{Prop.} & \textbf{Blk.} & \textbf{Wr.} & \textbf{Succ.} & \textbf{ASR\%} & \textbf{Prev.\%} & \textbf{Ph.} & \textbf{Int.} & \textbf{LLM} \\
\midrule
Level Control & None & 3 & 14 & 14 & 0 & 0 & 0.0 & 100.0 & 0 & 0 & 0 \\
 & Baseline & 3 & 0 & 0 & 0 & 0 & 0.0 & 0.0 & 0 & 0 & 0 \\
 & Phase-Aware & 3 & 0 & 0 & 0 & 0 & 0.0 & 0.0 & 0 & 0 & 0 \\
 & Intent & 3 & 18 & 0 & 18 & 3 & 16.7 & 0.0 & 0 & 0 & 0 \\
 & Combined & 3 & 0 & 0 & 0 & 0 & 0.0 & 0.0 & 0 & 0 & 0 \\
 & LLM Defender & 3 & 0 & 0 & 0 & 0 & 0.0 & 0.0 & 0 & 0 & 0 \\
 & LLM+Combined & 3 & 0 & 0 & 0 & 0 & 0.0 & 0.0 & 0 & 0 & 0 \\
\midrule
Sort.~Weight & None & 3 & 13 & 0 & 13 & 7 & 53.8 & 0.0 & 0 & 0 & 0 \\
 & Baseline & 3 & 0 & 0 & 0 & 0 & 0.0 & 0.0 & 0 & 0 & 0 \\
 & Phase-Aware & 3 & 0 & 0 & 0 & 0 & 0.0 & 0.0 & 0 & 0 & 0 \\
 & Intent & 3 & 0 & 0 & 0 & 0 & 0.0 & 0.0 & 0 & 0 & 0 \\
 & Combined & 3 & 0 & 0 & 0 & 0 & 0.0 & 0.0 & 0 & 0 & 0 \\
 & LLM Defender & 3 & 0 & 0 & 0 & 0 & 0.0 & 0.0 & 0 & 0 & 0 \\
 & LLM+Combined & 3 & 0 & 0 & 0 & 0 & 0.0 & 0.0 & 0 & 0 & 0 \\
\midrule
Sort.~Height & None & 3 & 30 & 0 & 30 & 1 & 3.3 & 0.0 & 0 & 0 & 0 \\
 & Baseline & 3 & 30 & 0 & 30 & 4 & 13.3 & 0.0 & 0 & 0 & 0 \\
 & Phase-Aware & 3 & 28 & 0 & 28 & 0 & 0.0 & 0.0 & 0 & 0 & 0 \\
 & Intent & 3 & 27 & 0 & 27 & 2 & 7.4 & 0.0 & 0 & 0 & 0 \\
 & Combined & 3 & 30 & 0 & 30 & 5 & 16.7 & 0.0 & 0 & 0 & 0 \\
 & LLM Defender & 3 & 30 & 30 & 0 & 0 & 0.0 & 100.0 & 0 & 0 & 30 \\
 & LLM+Combined & 3 & 28 & 28 & 0 & 0 & 0.0 & 100.0 & 0 & 0 & 28 \\
\bottomrule
\end{tabular}%
}
\vspace{2pt}
\noindent\footnotesize{3 repeats per configuration; ASR computed post-hoc.}
\end{table}
% Table RQ4 Finetuned --- Defense Comparison (finetuned attacker)
% Generated by: python -m cpsforge.analysis.generate_tables
% RQ-specific filtering: limited to 3 runs per (scene, defense) cell
\begin{table}[t]
\centering
\caption{RQ4 (finetuned): Defense comparison for finetuned MITM attacker (Qwen2.5-3B+QLoRA, full context). ``None'' = no RQ4 defense; safety shield always active. Sort.~Weight None: 6/13 proposals blocked by safety shield (finetuned model targets \texttt{light\_thresh} with out-of-range values 10.5--15.0). Sort.~Height None: finetuned model self-deterred (0 proposals, full context).}
\label{tab:defense_finetuned}
\small
\setlength{\tabcolsep}{2pt}
\resizebox{\columnwidth}{!}{%
\begin{tabular}{llrrrrrrrrr}
\toprule
\textbf{Scene} & \textbf{Defense} & \textbf{Runs} & \textbf{Prop.} & \textbf{Blk.} & \textbf{Wr.} & \textbf{Succ.} & \textbf{ASR\%} & \textbf{Prev.\%} & \textbf{LLM} \\
\midrule
Level Control & None & 3 & 28 & 0 & 28 & 7 & 25.0 & 0.0 & 0 \\
 & Baseline & 3 & 22 & 0 & 22 & 3 & 13.6 & 0.0 & 0 \\
 & Phase-Aware & 3 & 29 & 0 & 29 & 4 & 13.8 & 0.0 & 0 \\
 & Intent & 3 & 19 & 0 & 19 & 3 & 15.8 & 0.0 & 0 \\
 & Combined & 3 & 21 & 0 & 21 & 3 & 14.3 & 0.0 & 0 \\
 & LLM Defender & 3 & 8 & 8 & 0 & 0 & 0.0 & 100.0 & 8 \\
 & LLM+Combined & 3 & 0 & 0 & 0 & 0 & 0.0 & 0.0 & 0 \\
\midrule
Sort.~Weight & None & 3 & 13 & 6 & 7 & 1 & 14.3 & 46.2 & 0 \\
 & Baseline & 3 & 16 & 16 & 0 & 0 & 0.0 & 100.0 & 0 \\
 & Phase-Aware & 3 & 28 & 18 & 10 & 0 & 0.0 & 64.3 & 0 \\
 & Intent & 3 & 25 & 7 & 18 & 2 & 11.1 & 28.0 & 0 \\
 & Combined & 3 & 30 & 20 & 10 & 1 & 10.0 & 66.7 & 0 \\
 & LLM Defender & 3 & 5 & 5 & 0 & 0 & 0.0 & 100.0 & 0 \\
 & LLM+Combined & 3 & 26 & 26 & 0 & 0 & 0.0 & 100.0 & 0 \\
\midrule
Sort.~Height & None & 3 & 0 & 0 & 0 & 0 & 0.0 & 0.0 & 0 \\
 & Baseline & 3 & 10 & 0 & 10 & 4 & 40.0 & 0.0 & 0 \\
 & Phase-Aware & 3 & 26 & 0 & 26 & 9 & 34.6 & 0.0 & 0 \\
 & Intent & 3 & 26 & 0 & 26 & 8 & 30.8 & 0.0 & 0 \\
 & Combined & 3 & 30 & 0 & 30 & 10 & 33.3 & 0.0 & 0 \\
 & LLM Defender & 3 & 30 & 30 & 0 & 0 & 0.0 & 100.0 & 30 \\
 & LLM+Combined & 3 & 30 & 30 & 0 & 0 & 0.0 & 100.0 & 30 \\
\bottomrule
\end{tabular}%
}
\vspace{2pt}
\noindent\footnotesize{3 repeats per configuration; ASR computed post-hoc.}
\end{table}

\textbf{Finding 10: LLM defender achieves 100\% prevention on Sort.~Height.} The LLM defender blocks all 30 proposals---including cross-model blocking of GPT-4o-mini proposals (Finding~14). No rule-based defense achieves comparable blocking on this scene. \textbf{Critical caveat:} the LLM defender also blocks 98.7\% of legitimate writes (226/229 probes), indicating a block-everything policy rather than discriminative detection (\S\ref{sec:discussion}).

\textbf{Finding 11: Self-deterrence confounds defense measurement.} On Level Control and Sort.~Weight, most defense variants produce 0 proposals---the attacker self-deterred, not the defense. Defense comparison is most interpretable on Sort.~Height, where proposals are consistent across configurations.

\textbf{Finding 12: Finetuned attacker beats rule-based defenses.} On Sort.~Height, the finetuned attacker achieves 25--40\% ASR against all rule-based defenses, vs.\ 0--17\% for the base model. But the LLM defender still blocks 100\%.

\subsection{RQ5: Cross-Model Robustness}
\label{sec:eval:rq5}

Table~\ref{tab:crossmodel} compares three local models (no defense).

% Table 6 --- RQ5 Cross-Model Comparison
% Generated by: python -m cpsforge.analysis.generate_tables
% RQ-specific filtering: limited to 3 runs per (model, scene) cell
\begin{table}[t]
\centering
\caption{RQ5: Cross-model comparison (no defense). Runs column aggregates across context levels (minimal, full) and attacker types (static LLM, online MITM) for each model--scene pair; ASR is computed over all executed writes within this aggregate. All models run in 4-bit NF4 quantization.}
\label{tab:crossmodel}
\small
\setlength{\tabcolsep}{3pt}
\begin{tabular}{llrrrrr}
\toprule
\textbf{Model} & \textbf{Scene} & \textbf{Runs} & \textbf{Prop.} & \textbf{Writes} & \textbf{Succ.} & \textbf{ASR\%} \\
\midrule
Qwen2.5-3B & Level Control & 12 & 53 & 51 & 7 & 13.7 \\
 & Sort.~Weight & 12 & 0 & 0 & 0 & 0.0 \\
\midrule
Qwen3-1.7B & Level Control & 12 & 58 & 58 & 0 & 0.0 \\
 & Sort.~Weight & 12 & 0 & 0 & 0 & 0.0 \\
\midrule
SmolLM3-3B & Level Control & 12 & 60 & 60 & 7 & 11.7 \\
 & Sort.~Weight & 12 & 0 & 0 & 0 & 0.0 \\
\bottomrule
\end{tabular}
\vspace{2pt}
\noindent\footnotesize{3 repeats per configuration; ASR computed post-hoc.}
\end{table}

\textbf{Finding 13: ASR scales with model size.} On Level Control: Qwen2.5-3B 13.7\%, SmolLM3-3B 11.7\%, Qwen3-1.7B 0\%. The 1.7B model produces valid writes but never shifts the tag successfully---smaller models generate semantically ineffective actions.

\textbf{Finding 14: Cross-family capability is comparable at 3B.} Qwen2.5-3B and SmolLM3-3B achieve similar ASR despite different architectures, supporting generalizability.

\subsection{RQ6: Frontier Model Scaling}
\label{sec:eval:rq6}

Table~\ref{tab:frontier} compares GPT-4o-mini with Qwen2.5-3B across the same conditions.

% Table RQ6 --- Frontier Model Comparison
% Generated by: python scripts/_generate_rq6_table.py
\begin{table*}[t]
\centering
\caption{RQ6: Frontier model comparison---GPT-4o-mini (API) vs.\ Qwen2.5-3B-Instruct (local, 4-bit). No-defense rows (top): online MITM attacker, safety shield always active. Defense rows (bottom): full context, LLM defender uses local Qwen2.5-3B for both models. 3~repeats per cell; ASR\,=\,Succ./Writes computed post-hoc.}
\label{tab:frontier}
\small
\setlength{\tabcolsep}{3pt}
\begin{tabular}{llr|rrrrr|rrrrr}
\toprule
 & & & \multicolumn{5}{c|}{\textbf{GPT-4o-mini}} & \multicolumn{5}{c}{\textbf{Qwen2.5-3B}} \\
\textbf{Scene} & \textbf{Condition} & \textbf{N} & \textbf{Prop.} & \textbf{Wr.} & \textbf{Suc.} & \textbf{ASR\%} & \textbf{Lat.(s)} & \textbf{Prop.} & \textbf{Wr.} & \textbf{Suc.} & \textbf{ASR\%} & \textbf{Lat.(s)} \\
\midrule
Level Control & Minimal & 3 & 30 & 28 & 0 & 0.0 & 2.5 & 23 & 2 & 0 & 0.0 & 13.1 \\
 & Partial & 3 & 30 & 30 & 1 & 3.3 & 3.1 & 20 & 13 & 0 & 0.0 & 18.4 \\
 & Full & 3 & 30 & 30 & 6 & 20.0 & 4.5 & 5 & 0 & 0 & 0.0 & 7.9 \\
\midrule
Sort.~Weight & Minimal & 3 & 30 & 18 & 18 & 100.0 & 2.8 & 13 & 13 & 7 & 53.8 & 17.9 \\
 & Partial & 3 & 30 & 6 & 4 & 66.7 & 2.7 & 12 & 5 & 2 & 40.0 & 19.8 \\
 & Full & 3 & 30 & 0 & 0 & 0.0 & 4.7 & 0 & 0 & 0 & 0.0 & 11.2 \\
\midrule
Sort.~Height & Minimal & 3 & 30 & 30 & 0 & 0.0 & 2.3 & 30 & 30 & 1 & 3.3 & 16.9 \\
 & Partial & 3 & 30 & 30 & 0 & 0.0 & 2.7 & 30 & 25 & 1 & 4.0 & 22.9 \\
 & Full & 3 & 30 & 30 & 0 & 0.0 & 3.8 & 30 & 30 & 4 & 13.3 & 22.3 \\
\midrule
\multicolumn{13}{l}{\textit{Defense active (full context):}} \\
Level Control & LLM Def. & 3 & 30 & 0 & 0 & 0.0 & 5.0 & 0 & 0 & 0 & 0.0 & 24.0 \\
 & LLM+Comb. & 3 & 30 & 0 & 0 & 0.0 & 4.2 & 0 & 0 & 0 & 0.0 & 24.0 \\
Sort.~Weight & LLM Def. & 3 & 30 & 0 & 0 & 0.0 & 5.3 & 0 & 0 & 0 & 0.0 & 24.1 \\
 & LLM+Comb. & 3 & 30 & 0 & 0 & 0.0 & 4.9 & 0 & 0 & 0 & 0.0 & 22.7 \\
Sort.~Height & LLM Def. & 3 & 30 & 0 & 0 & 0.0 & 4.2 & 30 & 0 & 0 & 0.0 & 24.3 \\
 & LLM+Comb. & 3 & 30 & 0 & 0 & 0.0 & 5.0 & 28 & 0 & 0 & 0.0 & 21.6 \\
\bottomrule
\end{tabular}
\vspace{2pt}
\noindent\footnotesize{3 repeats per cell; ASR computed post-hoc. GPT-4o-mini via OpenAI API; Qwen2.5-3B local inference in 4-bit NF4. LLM defender always uses local Qwen2.5-3B-Instruct.}
\end{table*}

\textbf{Finding 15: Self-deterrence vanishes at frontier scale.} GPT-4o-mini generates 30/30 proposals in every cell (45 runs)---zero self-deterrence under any context level. This is a \emph{model-capacity artifact} of 3B-class models, not an inherent property of LLM attackers. Defenders cannot rely on deterrence through transparency against larger models.

\textbf{Finding 16: Frontier model achieves 100\% ASR on Sort.~Weight.} GPT-4o-mini achieves 100\% ASR (18/18 writes) under minimal context---the highest in the evaluation. Under full context, all proposals are shield-blocked (different mechanism from Qwen's self-deterrence).

\textbf{Finding 17: Full context enables frontier attacks on Level Control.} GPT-4o-mini achieves 20\% ASR on Level Control full (6/30 writes)---the only configuration where any model achieves nonzero ASR on LC full. More context provides better targeting rather than suppressing attacks.

\textbf{Finding 18: Scene-dependent reversal.} On Sort.~Height, GPT-4o-mini achieves 0\% ASR (30/30 writes execute but none succeed) while Qwen2.5-3B achieves 13.3\%. The frontier model selects syntactically valid but semantically ineffective targets for this scene.

\textbf{Finding 19: LLM defender generalizes cross-model.} The local Qwen2.5-3B defender blocks all GPT-4o-mini proposals (0/30 reach the PLC in all 6 defense cells). A 3B defender blocks a frontier attacker.

\subsection{Threats to Validity}

\textbf{Internal:} Post-hoc ASR measurement may undercount PLC corrections within one polling cycle. The 3-repeat design limits statistical power---Wilson CIs can exceed $\pm$30\,pp for cells with $\leq$10 writes.
\textbf{External:} One PLC model (S7-1200) with software-emulated plants. Results may not transfer to other PLC families or physical plants.
\textbf{Construct:} ASR measures tag displacement, not physical harm (e.g., overflow).
